\documentclass[11pt]{article}
\usepackage[margin=1in]{geometry}
\usepackage{amsmath,amssymb,amsfonts,mathtools}
\usepackage{array,booktabs,enumitem,graphicx,xcolor,hyperref}
\usepackage[T1]{fontenc}
\usepackage[utf8]{inputenc}
\title{The renormalization group flow in 2D N=2 SUSY Landau-Ginsburg models}
\author{Source-backed arXiv sample}
\date{}
\begin{document}
\maketitle
\begin{abstract}
We investigate the renormalization of N=2 SUSY L-G models with central charge \$c=3p/(2+p)\$ perturbed by an almost marginal chiral operator. We calculate the renormalization of the chiral fields up to \$gg\{\textasciicircum{}*\}\$ order and of nonchiral fields up to \$g(g\textasciicircum{}\{*\})\$ order. We propose a formulation of the nonrenormalization theorem and show that it holds in the lowest nontrivial order. It turns out that, in this approximation, the chiral fields can not get renormalized \$\textbackslash{}Phi\textasciicircum{}\{k\}=\textbackslash{}Phi\textasciicircum{}\{k\}\_\{0\}\$. The \$\textbackslash{}beta\$ function then remains unchanged \$\textbackslash{}beta=\textbackslash{}epsilon gr\$.
\end{abstract}
\section{References And Citations}
\label{sec:refs}
This sample cites source keys \cite{rg:fl1} and \cite{rg:fl2}. Section~\ref{sec:refs} and Equation~\ref{eq:source-demo} provide cross-reference coverage.
\begin{equation}
\label{eq:source-demo}
a^2 + b^2 = c^2
\end{equation}
\begin{thebibliography}{9}
\bibitem{rg:fl1} Source bibliography entry 1 extracted from arXiv metadata/source key `rg:fl1`.
\bibitem{rg:fl2} Source bibliography entry 2 extracted from arXiv metadata/source key `rg:fl2`.
\bibitem{rg:fl3} Source bibliography entry 3 extracted from arXiv metadata/source key `rg:fl3`.
\bibitem{rg:fl4} Source bibliography entry 4 extracted from arXiv metadata/source key `rg:fl4`.
\end{thebibliography}

\end{document}
